% !TeX root = ../navier-stokes-manuscript.tex
\subsection{Velocity, acceleration, jerk, and the Tower}\label{sub:BodyFiniteTimeResponseEscalation}

The Zeno picture compresses space and time together.  A smaller width is not
reached by naming a later scale; the velocity has to move through the change
during the shorter interval available.  Let \(X\) be a spatial Banach space in
which the classical responses are defined on every strict subinterval.  For
each \(n\geq0\),
\begin{equation}\label{eq:BodyResponseIncrement}
  \norm{\partial_t^nu(t_1)-\partial_t^nu(t_0)}_X
  \leq
  \int_{t_0}^{t_1}\norm{\partial_t^{n+1}u(s)}_X\,ds.
\end{equation}
If \(\partial_t^{n+1}u\) stayed uniformly bounded before a finite \(T_*\),
this inequality would keep \(\partial_t^nu\) uniformly bounded as well.
Reading that implication upward, any escape of \(u\) in \(X\) forces every
finite temporal response to lose uniform control in the same space:
\begin{equation}\label{eq:BodyConditionalTemporalEscape}
  \sup_{0<t<T_*}\norm{u(t)}_X=\infty
  \quad\Longrightarrow\quad
  \sup_{0<t<T_*}\norm{\partial_t^nu(t)}_X=\infty
  \quad\text{for every finite }n.
\end{equation}
This is a necessary consequence of a proposed escape, not a proof that the
escape occurs.  It tells us why the temporal responses have to be examined
together:
\[
  u,
  \qquad
  \partial_tu,
  \qquad
  \partial_t^2u,
  \qquad
  \partial_t^3u,
  \qquad
  \ldots
\]
Write
\begin{equation}\label{eq:TemporalTowerDefinitionBody}
  V_n:=\partial_t^nu,
  \qquad
  Q_n:=\partial_t^np,
  \qquad n\in\Nzero.
\end{equation}
These are the velocity, temporal acceleration, temporal jerk, and higher
Eulerian responses of the fluid.  The material acceleration is the full
combination
\[
  D_tu=\partial_tu+(u\cdot\nabla)u;
\]
the Tower follows \(\partial_t^nu\) at fixed spatial position and keeps the
transport term inside the recurrence below.
The Tower is not an extra model placed beside Navier--Stokes.  The original
equation generates every rung.

At the base,
\[
  V_1
  =\nu\Delta V_0-(V_0\cdot\nabla)V_0-\nabla Q_0.
\]
Differentiate once:
\[
  V_2
  =\nu\Delta V_1
  -(V_1\cdot\nabla)V_0
  -(V_0\cdot\nabla)V_1
  -\nabla Q_1.
\]
Differentiate again, and the two ways of placing one derivative on each
transport factor produce the middle coefficient:
\[
  V_3
  =\nu\Delta V_2
  -(V_2\cdot\nabla)V_0
  -2(V_1\cdot\nabla)V_1
  -(V_0\cdot\nabla)V_2
  -\nabla Q_2.
\]
The general row is the binomial expansion of the transport term:
\begin{equation}\label{eq:BodyTemporalVelocityRecurrence}
  V_{n+1}
  +\sum_{a=0}^{n}\binom{n}{a}(V_a\cdot\nabla)V_{n-a}
  +\nabla Q_n
  =\nu\Delta V_n,
  \qquad
  \nabla\cdot V_n=0.
\end{equation}

Pressure does not create a second Tower.  Taking the divergence of the base
equation and using the decaying whole-space normalization gives
\begin{equation}\label{eq:PressureNormalization}
  -\Delta p=\partial_i\partial_j(u_i u_j),
  \qquad
  p=R_iR_j(u_i u_j),
\end{equation}
where \(R_j=\partial_j(-\Delta)^{-1/2}\).  After \(n\) time derivatives,
\begin{equation}\label{eq:BodyTemporalPressureRecurrence}
  Q_n
  =R_iR_j
  \sum_{a=0}^{n}\binom{n}{a}(V_a)_i(V_{n-a})_j.
\end{equation}
The velocity rows determine \(Q_n\); its gradient returns to
\eqref{eq:BodyTemporalVelocityRecurrence} as part of the next response.  Every
rung retains transport, pressure, viscosity, and incompressibility.

The Tower is now fixed by the original equation: velocity generates
acceleration, acceleration generates jerk, and every later response retains
transport, pressure, incompressibility, and viscosity.  One question has been
left inside the temporal law.  Each rung contains \(\nu\Delta V_n\).  As the
responses climb in time, what does that Laplacian force us to carry in space?
